\chapter{人工智能}

\section{领域导言与演进背景}

\begin{preamblebox}[知识领域概述与变革]
人工智能 (AI)致力于研究与解决那些难以或无法通过传统算法有效求解的计算问题。这类问题往往需要模拟人类智能，其解决方案通常能够跨类别泛化至更广泛的同构问题。如今，人工智能技术在计算领域无处不在，深度赋能电子邮件、社交媒体、智能摄影、金融量化交易以及智能虚拟语音助理（如 Siri、Alexa 等）等日常生活应用。此外，这些技术也被广泛应用于自主智能体的设计与分析中，使其能够自主感知外部环境并进行理性交互，例如自动驾驶车辆与各类自主移动机器人。

从传统上看，人工智能融合了符号学派与亚符号/联结学派的方法。其求解机制依赖于一套兼具通用性与领域专用性的知识表示体系、问题求解机制与数值优化技术。这些技术涵盖了感知（如语音识别、自然语言理解、计算机视觉）、问题求解（如启发式搜索、自动规划、约束优化）、动作执行（如机器人学、任务自动化、智能控制），以及支撑上述能力的系统架构（如单智能体架构与多智能体系统）。学习人工智能旨在培养学生具备以下专业能力：针对特定现实问题准确研判是否适用 AI 方法；科学选择与设计适宜的知识表示与推理机制；高质量实现工程原型；并在算法计算性能与更广泛的社会影响两个维度上进行严谨的系统评估。

在过去十年中，在一系列高影响力的机器学习重大突破推动下，“人工智能”一词已在企业界、新闻媒体与日常话语中广泛普及。这些跨越式进展得益于海量大数据的可用性、计算算力（如 GPU/TPU）的迅猛飞跃以及算法理论的持续演进。特别显著的是，人工智能的研究范式发生了重大转向——从以往高度依赖人工工程化设计的特征与知识表示，转向通过大规模数据集优化自适应学习得到的特征表示。由此催生的技术变革使“神经网络”和“深度学习”成为大众熟知的流行语。企业纷纷将基于 AI 的技术方案作为核心服务增值亮点进行推广，使得“人工智能”既是一个严谨的技术学术术语，也是一个备受瞩目的商业营销关键词。与此同时，生物学、艺术创作、建筑设计和金融工程等跨学科领域也日益深度引入 AI 技术以求解其学科前沿问题。人类历史上首次迎来广泛大众能够直接接触并使用复杂精尖的 AI 驱动工具的新时代，包括根据提示词自动生成长篇论文或优美诗歌的生成式工具、依据文本描述生成高品质摄影或艺术图像的扩散工具，以及能够极其逼真地合成特定人物图像与音视频的深度伪造技术。AI 技术现已深度渗透至金融股票交易、新闻与社交网络信息流推荐算法、招聘求职者自动化多维筛选、临床医疗病症辅助诊断，乃至通过再犯预测算法深度介入司法量刑判决。因此，人工智能技术具备极为深远的社会影响力，在研发与落地应用的全生命周期中必须对其潜在风险进行审慎考量与科学治理。

\textbf{自 CS2013 以来的核心变革：}
\begin{enumerate}
    \item \textbf{领域命名规范化更替}：将 CS2013 中的“智能系统”正式更名为“人工智能”，以精准呼应学科内部最通用的学术规范及全社会最广泛的行业共识。
    \item \textbf{强化神经网络与表示学习，平衡符号与联结范式}：顺应深度学习与表示学习的技术浪潮，大幅提升了神经网络与特征学习的教学权重。图搜索作为 AI 的经典基石依然得到坚守，但适度精简了部分传统符号演算方法，以腾出学时深入剖析亚符号方法与学习型表征。与此同时，课程体系坚决保留基于知识的系统与符号推理，因为符号方法具有独到的可解释性与严密逻辑推理能力，在工程实践中依然发挥关键作用，且最新前沿“神经符号人工智能”正在将学习表征与符号知识深度融合。
    \item \textbf{突出跨领域工程落地与深度生成式模型}：全面深化 AI 在医疗健康、环境可持续发展、社交网络治理等多元垂直领域的实战应用。特别将以 ChatGPT、Bard/Gemini 为代表的大语言模型（LLM）以及以 DALL-E、Midjourney、Stable Diffusion 为代表的深度图像扩散生成模型纳入核心教学视野，系统讲授其宏观技术脉络、落地场景、幻觉缺陷与潜在风险。
    \item \textbf{全面融入社会影响、伦理道德与可解释性}：深刻体现理解并客观评估 AI 技术与应用对人类社会深远影响的重要性，将人工智能伦理、算法公平性、公众信任度与可解释性作为核心必修议题深度织入各知识模块。
    \item \textbf{构筑与数据科学及数据管理的紧密协同链接}：强化人工智能知识领域与数据管理（DM）及数据科学基础设施的跨学科交叉映射。
    \item \textbf{确立面向所有计算专业学生的人工智能素养与批判性思维目标}：鉴于 AI 技术已广泛渗透至计算机科学的各行各业，明确将培养每位计算机科学专业学生的基础 AI 素养与批判性思辨能力作为核心培养目标。
\end{enumerate}
\end{preamblebox}

\section{核心学时分配建议}

本知识领域在“向日葵模型”架构下的学时推荐如下表所示：

\begin{table}[H]
\centering
\setlength{\tabcolsep}{3.5pt}\fontsize{8pt}{9.6pt}\selectfont\begin{tabularx}{\columnwidth}{Y c c}
\toprule
\textbf{知识单元} & \textbf{计算机核心} & \textbf{领域核心} \\
\midrule
基础议题 & 2 & 1 \\
搜索算法与问题求解 & 2 (+3 AL)$^*$ & 4 \\
基础知识表示与推理 & 1 (+1 MSF)$^\dagger$ & 2 \\
机器学习 & 4 & 4 \\
应用与社会影响 & 2 & 2 \\
逻辑表示与推理 & -- & 选修 \\
概率表示与推理 & -- & 选修 \\
规划系统 & -- & 选修 \\
智能体系统 & -- & 选修 \\
自然语言处理 & -- & 选修 \\
机器人学 & -- & 选修 \\
感知与计算机视觉 & -- & 选修 \\
\midrule
\textbf{合计} & \textbf{11 (+4 跨领域共享)$^\ddagger$ 学时} & \textbf{13 学时} \\
\bottomrule
\end{tabularx}
\begin{flushleft}\footnotesize
$^*$注：搜索算法单元包含 5 个计算机核心学时，其中 3 个学时（无信息图搜索部分）与算法基础（算法/基础数据结构与算法）交叉共享。\\
$^\dagger$注：基础知识表示与推理单元包含 2 个计算机核心学时，其中 1 个学时（贝叶斯概率推理基础）与数学与统计学基础交叉共享。\\
$^\ddagger$注：计算机核心总学时为 11 学时（由 AI 独立承担），加上由 AL 和 MSF 共享承载的 4 学时，该领域覆盖的计算机核心内容总量共计 15 学时。
\end{flushleft}
\end{table}

\section{知识单元与详细大纲}

\subsection{人工智能基础议题}
\textbf{【学时建议】}：2 计算机核心学时，1 领域核心学时，包含选修主题。

本单元系统阐述人工智能的核心概念、发展全景与基本哲学命题，界定智能行为与智能体的基本内涵，为学生构建起涵盖问题建模、环境属性分类与社会影响洞察的全局认知视野。

\subsubsection{计算机核心主题}
\begin{itemize}
    \item \textbf{人工智能全貌与前沿进展}：人工智能求解问题的本质范式；近期标志性成功工程应用案例（如大语言模型、视觉生成模型、AlphaFold 蛋白质结构预测、高阶自动驾驶等）。
    \item \textbf{智能体的概念与典型范式}：智能体的形式化定义；反应式智能体与深思式智能体的核心特征与对比。
    \item \textbf{智能行为的本质思考}：什么是智能行为？智能行为的计算模拟途径。
    \item \textbf{图灵测试与心智思辨}：图灵测试的提出背景、测试协议及其理论缺陷；经典哲学反思与认知争议。
    \item \textbf{多模态交互}：多模态输入与输出（文本、视觉、语音、传感器时序等维度的融合感知与多模态生成）。
    \item \textbf{推理范式辨析}：理性推理与非理性/类人推理的本质边界与目标差异。
    \item \textbf{问题与环境特征分类体系}：
    \begin{itemize}
        \item 完全可观察 vs. 部分可观察；
        \item 单智能体 vs. 多智能体；
        \item 确定性 vs. 随机性；
        \item 静态 vs. 动态；
        \item 离散 vs. 连续。
    \end{itemize}
    \item \textbf{智能体本质与自主性特征}：
    \begin{itemize}
        \item 完全自主、半自主与混合主动自主；
        \item 条件反射式、基于目标与基于效用的智能体架构；
        \item 不确定性环境与不完全信息条件下的决策制定机制；
        \item 环境感知与动态物理交互的核心重要性；
        \item 基于学习的智能体演进机制；
        \item 具身智能体核心三要素：传感器、系统动力学与执行器。
    \end{itemize}
    \item \textbf{人工智能的应用扩张与社会冲击}：AI 技术的爆发式增长及其在经济结构、社会就业与伦理治理层面的广泛影响。
\end{itemize}

\subsubsection{领域核心主题}
\begin{itemize}
    \item \textbf{环境与问题特征深度剖析}：结合复杂现实工程实例（如高频量化交易、开放道路自动驾驶、智能电网调度），深入解析各类环境维度的复合交叉特征。
    \item \textbf{智能体系统架构进阶}：结合典型工业级系统，深入分析分层控制、意图驱动与效用函数在复杂智能体中的架构实现。
    \item \textbf{人工智能深远影响专题研讨}：全面评估 AI 产业应用带来的宏观经济效益、劳动力结构重塑、数据隐私边界与伦理合规挑战。
\end{itemize}

\subsubsection{选修主题}
\begin{itemize}
    \item \textbf{人工智能哲学与认知科学前沿}：强人工智能与弱人工智能的思辨；功能主义、机器意识与意向性理论；机器伦理学与机器权利哲学探讨。
\end{itemize}

\subsubsection{示范性学习成果与能力目标}
\begin{enumerate}
    \item \textbf{阐释}图灵测试的基本协议机制，并结合约翰·塞尔的“中文屋”思想实验，深入剖析形式符号操作与真正语义理解之间的本质鸿沟。
    \item \textbf{辨析}最优理性推理/行为与拟人化推理/行为的根本差异，能够依据特定业务需求确定系统设计应追求的目标范式。
    \item \textbf{界定}任意现实工程问题与应用场景在可观察性、随机性、动态性、离散/连续性以及智能体协作模式等维度上的核心环境特征。
\end{enumerate}

\subsection{搜索算法与问题求解}
\textbf{【学时建议】}：5 计算机核心学时（其中 3 学时与 AL 算法的无信息搜索交叉），4 领域核心学时，包含丰富选修主题。

本单元系统讲授人工智能中基于状态空间搜索的问题求解方法，涵盖盲目图搜索、启发式图搜索、博弈对抗搜索及现代局部/随机优化搜索技术。

\subsubsection{计算机核心主题}
\begin{itemize}
    \item \textbf{问题的状态空间表示}：
    \begin{itemize}
        \item 状态、初始状态、目标测试与动作算子的规范化形式定义；
        \item 状态的特征因子化表示与假设空间；
        \item 基于图搜索的问题求解本质：将图视作潜在的状态空间，将树遍历视作对该状态空间的系统性探索；
        \item 隐式状态图的动态展开机制（状态图并非预先显式完整存储，而是通过后继生成函数动态按需构建）。
    \end{itemize}
    \item \textbf{无信息图搜索技术}：
    \begin{itemize}
        \item 广度优先搜索 (BFS)及其在无权最优路径求解中的应用；
        \item 深度优先搜索 (DFS)及其内存优势；
        \item 迭代加深深度优先搜索 (IDS)的完备性与时空折中；
        \item 一致代价搜索（Uniform Cost Search, UCS / 广义 Dijkstra 算法）。
    \end{itemize}
    \item \textbf{启发式/有信息图搜索技术}：
    \begin{itemize}
        \item 启发函数的数学构建及其可采纳性与一致性；
        \item 局部搜索与爬山法；
        \item 局部极小值与搜索地形特征（搜索高地、山脊与全局最优解的辨析）；
        \item 贪婪最佳优先搜索；
        \item A* 搜索算法的核心原理、评价函数 $f(n) = g(n) + h(n)$ 及其最优性定理。
    \end{itemize}
    \item \textbf{搜索算法时空复杂度分析}：各经典搜索算法的时间复杂度与空间复杂度理论推导、完备性与最优性证明。
\end{itemize}

\subsubsection{领域核心主题}
\begin{itemize}
    \item \textbf{双向搜索}：正向搜索与逆向搜索的交汇原理与边界条件。
    \item \textbf{束搜索}：受限束宽贪婪剪枝搜索机制。
    \item \textbf{双人对抗博弈搜索}：
    \begin{itemize}
        \item 极小化极大搜索算法与递归博弈树展开；
        \item Alpha-Beta 剪枝技术（$\alpha$-$\beta$ Pruning）及其最高剪枝效率（最优走法排序下的渐近加速）；
        \item 搜索深度截断与启发式静态估值函数设计。
    \end{itemize}
    \item \textbf{A* 搜索算法的工程优化实现}：利用优先队列、开放列表与闭合列表优化大规模状态搜索性能。
\end{itemize}

\subsubsection{选修主题}
\begin{itemize}
    \item \textbf{深入洞察复杂搜索空间}：
    \begin{itemize}
        \item 搜索树构建原理与分支因子导致的组合爆炸；
        \item 动态不确定搜索空间建模；
        \item 复杂搜索地形拓扑（鞍点 Saddle Points、高维平坦高原 Plateau、狭窄山脊 Ridges 等）。
    \end{itemize}
    \item \textbf{约束满足问题 (CSP)}：约束图表示、回溯搜索、前向检验与弧相容算法（AC-3）。
    \item \textbf{禁忌搜索}：短时禁忌表与特赦准则。
    \item \textbf{内存受限与高阶 A* 变体}：IDA*（迭代加深 A*）、SMA*（简化内存受限 A*）、RBFS（递归最佳优先搜索）。
    \item \textbf{双人博弈高级专题}：地平线效应与静态搜索；开局库与残局完备解库；博弈弱解与强解的数学定义（如西洋跳棋 Checkers 已被完全弱解）。
    \item \textbf{算法实战优化}：极小化极大算法与束搜索的高性能工程实现。
    \item \textbf{期望极大搜索}：含随机概率机会节点的博弈树展开与马尔可夫决策过程求解初阶。
    \item \textbf{随机化与演化搜索算法}：
    \begin{itemize}
        \item 模拟退火算法与 Metropolis 接受准则、降温调度表设计；
        \item 遗传算法：染色体编码、轮盘赌选择、交叉、变异与适应度函数；
        \item 蒙特卡洛树搜索 (MCTS)：选择、扩展、模拟与反向传播（Backpropagation / UCB1 策略）。
    \end{itemize}
\end{itemize}

\subsubsection{示范性学习成果与能力目标}
\begin{enumerate}
    \item \textbf{设计}经典组合谜题（如 N 皇后问题、滑块拼图或水桶分水问题）的形式化状态空间表示方案。
    \item \textbf{选择并实现}适宜的无信息搜索算法（如 BFS、DFS、IDS）求解具体问题（如迷宫遍历或井字棋），并严格推导其时间与空间复杂度。
    \item \textbf{设计}具备启发性的估值函数，并在特定问题场景下（如移动机器人在二维网格地图中的避障路径规划）\textbf{编码实现}有信息启发式搜索。
    \item \textbf{论证并评估}给定启发函数是否满足可采纳性与一致性，判断其能否严格保证搜索得到最优解。
    \item \textbf{ 应用}极小化极大搜索算法求解双人零和对抗博弈（如四子棋 Connect Four），在达到指定搜索深度截断时调用静态启发估值函数进行分数回溯传递。
    \item \textbf{设计并实现}遗传算法，求解复杂的组合优化或排班调度问题。
    \item \textbf{设计并实现}模拟退火冷却计划表，帮助局部搜索过程有效跳出局部极小值陷阱。
    \item \textbf{设计并编码实现} A* 搜索与束搜索算法求解现实工程问题，并在解的最优性、扩展节点数及运行开销等维度上与其他搜索算法进行定量基准评测。
    \item \textbf{应用} Alpha-Beta 剪枝技术大幅削减双人博弈树的搜索节点规模，验证其与朴素 Minimax 搜索结果的一致性。
    \item \textbf{对比分析}遗传算法与经典基于树/图的搜索策略及数值梯度优化算法（如梯度下降法）的异同，阐明在何种特定问题场景下应优先采用演化算法。
    \item \textbf{综合权衡}各类启发式搜索策略在面对不同规模、复杂地形与算力约束问题时的有效性与工程适用度。
\end{enumerate}

\subsection{基础知识表示与推理}
\textbf{【学时建议】}：2 计算机核心学时（其中 1 学时与 MSF 离散概率交叉），2 领域核心学时。

本单元系统讲授知识表示的核心范式，从传统形式逻辑符号表示延伸至现代连续亚符号表征，并以贝叶斯理论为基础构建不确定性知识推理与理性决策的坚实理论底座。

\subsubsection{计算机核心主题}
\begin{itemize}
    \item \textbf{知识表示的核心类型}：
    \begin{itemize}
        \item 符号化与逻辑表示：将自然语言业务问题准确形式化为符号逻辑陈述；
        \item 学习驱动的亚符号表示：低维稠密向量与连续嵌入表示；
        \item 概率图模型初阶：朴素贝叶斯模型与贝叶斯信念网络直观拓扑。
    \end{itemize}
    \item \textbf{概率推理回顾与贝叶斯定理}：先验概率、后验概率、似然度；贝叶斯定理（Bayes' Theorem）的形式化陈述（与 MSF 离散概率交叉）。
    \item \textbf{贝叶斯推理基础}：基于观测证据进行假设后验概率计算与置信度更新。
\end{itemize}

\subsubsection{领域核心主题}
\begin{itemize}
    \item \textbf{随机变量与概率分布}：离散/连续随机变量、联合概率分布、边缘概率分布与条件概率分布。
    \item \textbf{概率公理体系}：柯尔莫哥洛夫公理及推论。
    \item \textbf{概率推断机制}：基于联合概率分布全表展开推断及其指数级复杂度瓶颈。
    \item \textbf{贝叶斯法则的数学推导}：从条件概率定义严谨推导贝叶斯法则。
    \item \textbf{复杂情境下的贝叶斯推理}：多证据源、序贯观测下的后验概率链式更新计算。
    \item \textbf{独立性与条件独立性}：边际独立性定义；条件独立性的数学定义及其对因果解耦的核心价值。
    \item \textbf{效用理论与理性决策}：偏好关系、效用函数与最大期望效用 (MEU)决策准则。
\end{itemize}

\subsubsection{示范性学习成果与能力目标}
\begin{enumerate}
    \item 给定自然语言描述的实际工程业务需求，将其准确\textbf{形式化编码}为符号逻辑或结构化知识表示。
    \item 结合贝叶斯定理与效用理论，深入\textbf{阐明}智能体如何在不确定性与信息不完全环境下做出最优理性决策。
    \item 运用贝叶斯定理在现实工程场景（如垃圾邮件识别、疾病筛查检测）中\textbf{执行概率推断}，准确推算在观测到特定证据条件下的假设后验概率。
    \item 熟练\textbf{应用}贝叶斯法则推导并求解多条件下的逆向概率推断问题。
    \item 准确\textbf{计算}各类随机事件的概率分布，并定量严格\textbf{检验}变量或事件之间是否满足统计独立性或条件独立性。
\end{enumerate}

\subsection{机器学习}
\textbf{【学时建议】}：4 计算机核心学时，4 领域核心学时，包含丰富选修主题。

本单元系统讲授机器学习的核心范式、基本理论与工程全流程，涵盖经典统计学习、深度神经网络、模型泛化评估以及数据偏见与算法伦理，培养学生构建工业级端到端预测模型与严谨评测系统的专业能力。

\subsubsection{计算机核心主题}
\begin{itemize}
    \item \textbf{机器学习任务的核心范式与定义}：
    \begin{itemize}
        \item 监督学习：分类（Classification，离散标签预测）与回归（Regression，连续数值预测）；
        \item 强化学习：智能体与环境通过“状态-动作-奖励”闭环试错交互学习最优策略；
        \item 无监督学习：聚类分析、数据降维与密度估计。
    \end{itemize}
    \item \textbf{机器学习基本哲学与理论基石}：
    \begin{itemize}
        \item 没有免费的午餐定理及其对算法普适性的启示；
        \item 不可判定性与学习理论的固有边界；
        \item 机器学习系统预测误差的本质来源（固有噪声、模型偏差与数据方差）。
    \end{itemize}
    \item \textbf{经典统计监督学习模型初阶}：
    \begin{itemize}
        \item 朴素贝叶斯分类器与决策树；
        \item 重点聚焦于直观机制与工程调用，不陷入底层繁复的数值优化推导，确保学生能准确理解并正确使用工业级开源实现。
    \end{itemize}
    \item \textbf{过拟合问题与模型复杂度调控}：
    \begin{itemize}
        \item 过拟合与欠拟合的现象及成因；
        \item 正则化与剪枝的直观调控逻辑；
        \item 偏差-方差权衡的核心内涵。
    \end{itemize}
    \item \textbf{真实数据工程与数据处理}：
    \begin{itemize}
        \item 数据清洗与预处理的工程重要性及常见踩坑陷阱；
        \item 缺失值处理策略：统计插补补全与显式标记缺失的技术差异及其对下游模型行为的深远影响；
        \item 类别变量编码（One-Hot 独热编码、标签序数编码）与连续实值特征变换；
        \item 数据归一化与标准化（Min-Max 缩放与 Z-score 标准化）；
        \item 坚决强调面对真实世界脏数据的处理实战，而非理想化的人工教材玩具数据。
    \end{itemize}
    \item \textbf{特征与假设表示}：
    \begin{itemize}
        \item 假设空间及其容量复杂度；
        \item 简单的基特征扩展（Basis Feature Expansion，如单变量二次方项特征构造）；
        \item 自动学习得到的特征表示（Learned Feature Representations / 特征表示学习）。
    \end{itemize}
    \item \textbf{严谨的机器学习评估方法论}：
    \begin{itemize}
        \item 训练集、验证集与独立测试集的严格隔离规范与防数据泄漏机制；
        \item 分类器常用评估指标体系（准确率 Accuracy、精确率 Precision、召回率 Recall、F1 分数等）；
        \item 基于未参与训练的保留测试集进行泛化性能无偏评估；
        \item 基于验证集进行模型超参数调优；
        \item 深入理解模型真实决策机制、明晰模型适用边界/盲区漏洞及其现实决策风险的关键意义。
    \end{itemize}
    \item \textbf{基础神经网络导论}：
    \begin{itemize}
        \item 人工神经元、多层感知机（MLP）拓扑结构、激活函数的非线性映射；
        \item 前向传播与反向传播模型训练机制的直观理解，不展开繁复矩阵微积分推导。
    \end{itemize}
    \item \textbf{机器学习工程伦理与偏见防范}：
    \begin{itemize}
        \item 基于真实商业场景与典型社会案例研究；
        \item 数据集偏见、算法设计偏见与评估基准偏见的识别。
    \end{itemize}
\end{itemize}

\subsubsection{领域核心主题}
\begin{itemize}
    \item \textbf{机器学习形式化优化建模}：
    \begin{itemize}
        \item 以最小二乘线性回归与对数几率回归为例建立优化数学模型；
        \item 目标损失函数的构建；
        \item 梯度下降法与随机梯度下降（SGD）的迭代更新机理；
        \item 抑制过拟合的正则化数学公式表达（$L_2$ 岭回归与 $L_1$ Lasso 正则项）。
    \end{itemize}
    \item \textbf{模型集成机制}：简单的加权多数表决法。
    \item \textbf{深度学习进阶}：
    \begin{itemize}
        \item 深度前馈神经网络的深层表达优势（直观机理分析）；
        \item 卷积神经网络 (CNN)的局部感受野、权重共享与池化下采样机理；
        \item 深度网络各层自动学习学到的特征表征的可视化分析。
    \end{itemize}
    \item \textbf{高阶性能评估与验证}：
    \begin{itemize}
        \item 分类高级指标（精确率、召回率、PR 曲线、Macro/Micro 平均）；
        \item 回归模型评估指标（均方误差 MSE、均方根误差 RMSE、平均绝对误差 MAE、$R^2$ 判定系数）；
        \item 混淆矩阵的构建与诊断；
        \item 交叉验证（Cross-Validation，如 $k$ 折交叉验证）；
        \item 结合交叉验证的网格搜索与随机搜索超参数调优。
    \end{itemize}
    \item \textbf{强化学习全景概览}：状态空间、动作空间、即时奖励、折扣因子、价值函数与 Q 学习直观机理。
    \item \textbf{机器学习重大跨行业工程落地剖析}：深入解析至少两个真实行业落地案例（如智慧医疗疾病诊断、量化金融风控、机器视觉质检、自然语言文本摘要、机器人运动控制、棋盘游戏博弈等）。
    \item \textbf{机器学习伦理深化}：真实场景案例研究；用户隐私保护（差分隐私初步）；算法公平性量化度量（统计均等、机会均等等算法防歧视准则）。
\end{itemize}

\subsubsection{选修主题}
\begin{itemize}
    \item \textbf{统计学习理论与参数估计}：极大似然估计（MLE）、最大后验估计（MAP）、期望最大化（EM）算法。
    \item \textbf{经典监督学习算法族深入探究}：
    \begin{itemize}
        \item 决策树划分准则（信息增益、信息增益率、基尼系数）与后剪枝算法；
        \item $k$ 近邻算法（$k$-NN）及其高维距离度量；
        \item 多层感知机反向传播链式求导法则的严密数学推导；
        \item 线性回归正规方程解析解与数值优化解对比；
        \item 逻辑回归似然函数求解与几率比分析；
        \item 支持向量机（SVM）理论：最大间隔超平面、对偶问题转换、支持向量判定与核技巧（高斯 RBF 核、多项式核）；
        \item 高斯过程非参数贝叶斯建模与核协方差函数。
    \end{itemize}
    \item \textbf{过拟合与高维学习理论}：维数灾难；$L_1$ 产生稀疏解与 $L_2$ 权重平滑收缩的严密几何与优化解析。
    \item \textbf{科学实验设计与高阶评测}：
    \begin{itemize}
        \item 分层抽样、时序分割与特征穿越防御；
        \item 受试者工作特征曲线（ROC 曲线）与曲线下面积（AUC）数学内涵。
    \end{itemize}
    \item \textbf{贝叶斯学习高阶}：朴素贝叶斯与广义线性模型的代数等价性；贝叶斯信念网络结构学习与参数估计；共轭先验分布与后验更新。
    \item \textbf{生成式模型理论}。
    \item \textbf{深度学习前沿架构与理论}：
    \begin{itemize}
        \item 深度前馈网络、神经正切核 (NTK)理论与过参数化网络训练动力学；
        \item 现代 CNN 经典拓扑（残差连接 ResNet、特征金字塔 FPN）；
        \item 自编码器与变分自编码器（VAE）；
        \item 循环神经网络（RNN、LSTM、GRU）与时间反向传播（BPTT）；
        \item 特征表示迁移、预训练微调与领域自适应；
        \item 对抗训练与生成对抗网络（GAN）。
    \end{itemize}
    \item \textbf{高阶特征表示工程}：手工特征提取算子；高维核基函数展开；深度自监督表征学习。
    \item \textbf{无监督聚类前沿}：$K$-均值聚类（$K$-Means）及其 $K$-Means++ 初始化；高斯混合模型（GMM）；自组织特征映射网络（SOM）。
    \item \textbf{半监督学习与弱监督学习}：伪标签技术与一致性正则化。
    \item \textbf{概率图模型进阶}：马尔可夫随机场（MRF）与因果推断图。
    \item \textbf{集成学习高级范式}：自举汇聚法、提升法（Boosting / AdaBoost、GBDT、XGBoost、LightGBM）、随机森林与门控专家混合模型。
    \item \textbf{统计学习理论极限}：为什么机器能够实现泛化；VC 维（Vapnik-Chervonenkis Dimension）；Rademacher 复杂度；概率近似正确（PAC）理论框架与泛化误差界。
    \item \textbf{强化学习前沿理论}：
    \begin{itemize}
        \item 探索与利用权衡（Exploration vs. Exploitation Trade-off，多臂老虎机模型）；
        \item 马尔可夫决策过程（MDP）严格数学规范与贝尔曼最优方程；
        \item 动态规划求解：价值迭代与策略迭代；
        \item 无模型强化学习：时间差分学习 TD($\lambda$)、Q-Learning、SARSA；
        \item 策略梯度方法（Policy Gradient Methods / REINFORCE 算法）；
        \item 深度强化学习（Deep Q-Networks / DQN、Actor-Critic 架构、PPO 近端策略优化算法）。
    \end{itemize}
    \item \textbf{推荐系统架构}：基于用户的协同过滤、矩阵分解（SVD）、图神经推荐算法。
    \item \textbf{机器学习专用计算硬件}：GPU 体系结构与 CUDA 并行加速、TPU 张量处理器、混合精度训练与分布式数据并行（DDP）。
    \item \textbf{多行业实战专题}：医学影像诊断、量化经济分析、智能辅导教育、计算机视觉目标检测、机器翻译 NLP、具身机器人、复杂博弈对抗。
    \item \textbf{数据挖掘深度交叉}：关联规则挖掘（Apriori 算法）、高维时序序列模式挖掘。
    \item \textbf{可信与可解释性人工智能 (Trustworthy \& Explainable AI, XAI)}：
    \begin{itemize}
        \item 数据集、算法与评估偏见的深度数学审计；
        \item 差分隐私联邦学习安全保护框架；
        \item 模型可解释性分析方法（SHAP 归因值、LIME 局部可解释替代模型、注意力权重热力图）。
    \end{itemize}
\end{itemize}

\subsubsection{示范性学习成果与能力目标}
\begin{enumerate}
    \item 深入\textbf{阐释}监督学习、强化学习与无监督学习三大核心学习范式在数据形式、优化目标与交互反馈上的根本区别。
    \item 准确\textbf{辨析}人工智能（AI）、机器学习（ML）与深度学习（DL）的层级从属关系与概念内涵。
    \item 将一个真实业务应用（如从手写像素网格图像中识别字母字符）\textbf{形式化建模}为规范的分类问题，明确界定可用输入特征维度与目标输出标签。
    \item \textbf{应用}两种或以上的经典统计学习算法（如 $k$ 近邻与逻辑回归）求解特定分类任务，并定量评估分类器的准确率与误差。
    \item 结合训练与验证学习曲线准确\textbf{识别}模型是否存在过拟合现象，并提出针对性的正则化、早停或剪枝等解决策略。
    \item \textbf{解释}机器学习算法如何在高维参数假设空间中转化为数学损失函数的搜索与连续优化迭代过程。
    \item 独立\textbf{编码实现}一个统计学习算法及其配套的梯度优化求解流程，完成模型参数训练，并在全新不可见测试样本上获得预测结果。
    \item \textbf{阐明}深度神经网络的前向计算与反向传播训练全过程，解释层级化特征表征的自动抽象机制。
    \item 深入\textbf{解释}规范严密的机器学习实验评测流程，辨析训练误差与泛化测试误差的根本差异，指出导致评测性能失真虚高的典型实验错误（如数据泄漏）。
    \item 针对基准数据集，从零\textbf{编码实现}包含数据清洗、特征预处理与无偏评估的方法管线，客观\textbf{对比}两种机器学习算法的泛化效能。
    \item 运用主流开源可视化套件（如 TensorBoard、WandB）\textbf{实时可视化}神经网络的训练损失曲线，并对网络隐藏层学到的特征表征进行降维可视化展示。
    \item 独立\textbf{编码实现}涵盖监督学习、强化学习与无监督学习的基础代表性算法原型。
    \item 针对特定的现实垂直领域问题，科学研判并\textbf{选择}最适宜的机器学习学习范式与技术路线。
    \item 深入\textbf{对比剖析}决策树、逻辑回归、朴素贝叶斯、多层神经网络与贝叶斯信念网络，结合典型应用案例说明各类算法的相对优势与最佳应用场景。
    \item 在包含缺失与噪声的复杂真实世界数据集上，独立部署并客观\textbf{评估}一个端到端预测系统的综合性能指标。
    \item 准确\textbf{概括}现代统计学习理论的学术前沿现状，客观评价其理论突破（如泛化误差界）及其在过参数化现代深度网络面前的理论局限。
    \item 深入\textbf{阐明}过拟合的发生机理，熟练掌握并\textbf{应用}交叉验证、早停、权重衰减（$L_2$ 正则化）与 Dropout 等综合防治策略。
    \item 深入\textbf{解析}机器学习中“假设空间容量”、“训练集样本规模”与“模型泛化精度”之间的三维制衡关系。
\end{enumerate}

\subsection{人工智能应用与社会影响}
\textbf{【学时建议】}：2 计算机核心学时，2 领域核心学时（与 SEP 交叉）。

\begin{infobox}[重要教学整合说明]
将人工智能的实际应用与伦理道德、算法公平性、公众信任度及可解释性议题，与具体的算法模型和数学理论有机交织、协同讲授，其教学成效远优于将伦理剥离为一门孤立专设的独立课程。在课程实施中，应尽可能将这些议题无缝嵌入到各技术单元中。例如，在讲授决策树算法时，可引导学生探讨如何将其应用于国土空间规划与土地资源可持续分配这一现实挑战，进而深入研讨自动化决策背后的环境生态公平性与伦理道德影响。

在核心学时分配上：\textbf{计算机核心} 必须至少深度覆盖一个具体工业级现实应用，并对 人工智能/ML 带来的广泛社会伦理议题建立系统概览；\textbf{领域核心} 则需进一步纵深拓展至一个或多个附加重大应用领域，深入剖析前沿深度生成式模型的工作机理，并对复杂的社会治理难题开展批判性剖析与研讨。
\end{infobox}

\subsubsection{核心主题大纲}
\begin{itemize}
    \item \textbf{人工智能在多元重大跨学科领域的实战赋能}：
    \begin{itemize}
        \item 深度覆盖智慧医疗与生命健康、环境与可持续计算、社交网络内容治理、数字金融与量化经济、智慧教育、特种机器人等广泛场景（CS 核心要求精讲至少 1 个场景，KA 核心要求拓展至少 1 个额外深度场景）；
        \item 将特定现实业务需求规范抽象并定义为 AI 形式化求解问题；
        \item 现实业务数据的可获取性、稀疏性与纯净度评估；
        \item 基础数据清洗与特征工程预处理准则；
        \item 数据集固有偏见：历史标注偏差、采样代表性失真；
        \item 算法模型偏见：目标函数损失惩罚倾斜导致的结构性失衡；
        \item 评测基准偏见：评测数据集与实际分布脱节导致的性能假象。
    \end{itemize}
    \item \textbf{前沿深度生成式模型}：
    \begin{itemize}
        \item 深度图像生成扩散模型（如截至 2023 年以来的 DALL-E、Midjourney、Stable Diffusion 等）：高层生成扩散原理、工业级商业赋能场景、版权争议、伪造滥用风险及其生成失效短板；
        \item 大语言模型（Large Language Models, LLMs，如截至 2023 年以来的 ChatGPT、Bard/Gemini 等）：自回归预训练与指令微调原理、工业级应用落地、事实性幻觉、安全对齐缺陷与逻辑推理脆弱性。
    \end{itemize}
    \item \textbf{人工智能深远社会影响与多维治理}：
    \begin{itemize}
        \item 人工智能伦理：算法决策责任归属、自主武器系统的伦理禁区、人类自主性与尊严捍卫；
        \item 算法公平性：不同群体间统计均等、均等化赔率等算法公平性数学指标，招聘与司法再犯预测中的偏见规避；
        \item 信任度与可解释性：深度黑盒模型的信任危机、决策可解释性（XAI）技术、全流程透明度与可审计性。
    \end{itemize}
\end{itemize}

\subsubsection{示范性学习成果与能力目标}
\begin{enumerate}
    \item 给定具体垂直行业的真实业务需求与复杂痛点，设计出规范严密的人工智能端到端解决方案，完整\textbf{明确}所需输入数据、清洗预处理管线、特征表示机制、核心 AI 算法选型以及全流程评估指标体系。
    \item 针对一个或多个在现实社会中已广泛部署落地的 AI 系统（如自动招聘初筛系统、智能信贷审批、人脸安防或司法量刑辅助系统），深入\textbf{分析}其广泛社会影响，精准识别并撰写关于伦理底线、算法公平性、数据偏见、用户信任与可解释性维度的专业风险评估报告。
    \item 深入\textbf{描述}当前主流文本大语言模型或图像扩散生成模型的核心失效模式（如事实幻觉、常识推理断裂、版权侵权与刻板偏见放大），客观\textbf{论述}这些局限性如何制约并影响其在高安全、高可靠关键任务场景中的工程应用与边界把控。
\end{enumerate}

\subsection{逻辑表示与推理 [选修]}
\textbf{【学时建议】}：拓展选修主题。

本单元系统涵盖基于经典逻辑与非经典逻辑的确定性知识表示与形式化推理体系，讲授本体工程、语义网以及常识推理和动作规划的形式化逻辑基础。

\subsubsection{选修主题详细大纲}
\begin{itemize}
    \item \textbf{命题逻辑与谓词逻辑回顾}：语法、语义模型、可满足性 SAT 与有效性（与 MSF 离散数学逻辑部分交叉）。
    \item \textbf{归结原理与定理自动证明}：
    \begin{itemize}
        \item 命题逻辑子句化与归结反演；
        \item 前向链接与后向链接推理算法。
    \end{itemize}
    \item \textbf{知识表示的核心工程难题}：表示充分性、推理效率与表达能力的平衡。
    \item \textbf{描述逻辑}：概念、角色、TBox（术语公理）与 ABox（断言公理）。
    \item \textbf{本体工程}：领域本体的设计原则、概念层级构建方法学。
    \item \textbf{语义网技术栈}：资源描述框架（RDF）、Web 本体语言（OWL）与大规模知识图谱。
    \item \textbf{非单调推理}：非经典逻辑、缺省逻辑与封闭世界假说 (CWA)。
    \item \textbf{形式化论辩理论}：抽象论辩框架与攻击关系。
    \item \textbf{关于动作与变化的形式化推理}：情景演算、事件演算与经典框架问题。
    \item \textbf{时间与空间推理}：阿兰时间区间代数（Allen's Interval Algebra）、定性空间拓扑关系。
    \item \textbf{逻辑程序设计}：Prolog 语言推理机制、答案集编程 (ASP)。
    \item \textbf{基于规则的专家系统}：产生式系统与 Rete 匹配算法。
    \item \textbf{语义网络} 与概念图。
    \item \textbf{基于模型的推理与基于案例的推理 (Model-based \& Case-based Reasoning, CBR)}。
\end{itemize}

\subsubsection{示范性学习成果与能力目标}
\begin{enumerate}
    \item 将自然语言表述的复杂业务陈述准确\textbf{翻译转化}为严格的一阶谓词逻辑公式。
    \item 将任意给定的命题逻辑或一阶逻辑公式正确\textbf{转换}为规范的合取范式子句集。
    \item 熟练\textbf{应用}归结反演原理对给定的一组逻辑子句进行机器推导，以形式化回答目标查询。
    \item 深入\textbf{对比剖析}结构化知识表示的主流模型（如框架系统、语义网络、本体论与描述逻辑），全面评述各自在知识表达能力与计算开销上的优势与局限。
    \item 深刻\textbf{阐明}非单调推理的核心要素，解释其作为人类常识信念系统与动态知识库维护机制的必要性。
    \item \textbf{对比分析}处理现实不确定性知识的各类形式化逻辑与定性技术手段。
    \item \textbf{对比分析}空间拓扑与定性物理关系的经典知识表示方法。
    \item \textbf{应用}情景演算或事件演算，对包含动态物理动作执行与环境状态变迁的复杂系统进行严密的形式化建模。
    \item 深入\textbf{阐释}时间推理与空间推理的本质特征、代数结构以及两者的时空协同交织关系。
    \item 深入\textbf{阐明}基于规则的推理、基于案例的推理（CBR）以及基于系统深层因果模型的推理机制的差异与应用场景。
    \item 给出形式化规划系统的严谨\textbf{定义}，阐明规划问题求解与经典状态图启发式搜索的本质区别。
    \item 深入\textbf{描述}基于状态搜索的规划、基于算子的规划与命题规划的技术路线，指明各自适用的任务环境。
    \item 深入\textbf{解释}单调推理与非单调推理的形式化边界，阐明新证据的引入在何种条件下会导致既往推论失效。
\end{enumerate}

\subsection{概率表示与推理 [选修]}
\textbf{【学时建议】}：拓展选修主题。

本单元系统讲授在不确定性、含噪声与不完全信息环境下构建概率知识表示与执行复杂概率推断的数学理论与工程实现。

\subsubsection{选修主题详细大纲}
\begin{itemize}
    \item \textbf{条件独立性回顾}：条件独立性代数公理化表述与图拓扑含义。
    \item \textbf{不确定性知识表示形式}：
    \begin{itemize}
        \item 贝叶斯网络：有向无环图结构与条件概率表（CPT）紧凑表示；
        \item 精确概率推断及其计算复杂度：变量消去法、团树算法以及精确推断的 NP-难性证明；
        \item 马尔可夫毯与有向分离定理（$d$-Separation）；
        \item 随机采样与近似推断：直接采样、拒绝采样、似然加权采样、马尔可夫链蒙特卡洛算法（MCMC，如 Gibbs 采样）；
        \item 马尔可夫网络（Markov Networks / 马尔可夫随机场 MRF）：无向图模型、势函数与吉布斯分布；
        \item 关系概率模型（Relational Probability Models / 统计关系学习）；
        \item 隐马尔可夫模型 (HMM)：前向-后向算法、Viterbi 最优路径解码算法与 Baum-Welch 参数学习算法。
    \end{itemize}
    \item \textbf{决策论高级专题}：
    \begin{itemize}
        \item 风险偏好公理与效用函数构造；
        \item 最大期望效用原则 (MEU)；
        \item 影响图（Influence Diagrams / 决策网络）。
    \end{itemize}
\end{itemize}

\subsubsection{示范性学习成果与能力目标}
\begin{enumerate}
    \item 给定具体贝叶斯网络的拓扑结构与条件概率分布参数，在观测到特定证据变量的条件下，准确\textbf{计算}待测假设变量的后验边缘概率分布。
    \item 深入\textbf{阐明}条件独立性断言如何切断图节点的稠密全连接，从而将全联合概率分布的指数级存储与计算开销急剧降低。
    \item 针对具体的复杂不确定性推理场景，举例并\textbf{构建}相应的概率图知识表示模型。
    \item 严格陈述贝叶斯网络精确概率推断的最坏情况计算复杂度，并系统\textbf{列举并对比}克服计算瓶颈的主流蒙特卡洛近似推断方法。
    \item 针对具体的现实不确定性决策场景，\textbf{设计并编程实现}至少一种概率图知识表示与推理求解引擎。
    \item 深入\textbf{剖析}时序概率推理系统（如动态贝叶斯网络与滤波模型）随时间推移面临的状态爆炸与推断复杂度挑战。
    \item \textbf{设计并编程实现}一个完整的隐马尔可夫模型（HMM），实现时序滤波估计或运用 Viterbi 算法解码最可能的隐藏时序状态序列。
    \item 深入\textbf{阐述}决策主体的主观偏好排序与实值效用函数之间的严密映射公理与数理逻辑关系。
    \item 深入\textbf{解释}如何将效用理论与概率图模型推理有机结合，以驱动自主智能体在不确定性环境下做出最大化期望效用的理性决策。
\end{enumerate}

\subsection{规划系统 [选修]}
\textbf{【学时建议】}：拓展选修主题。

本单元系统讲授自动化规划的核心理论，涵盖经典状态规划、偏序规划、分层任务网络（HTN）以及动态不确定环境下的闭环重规划机制。

\subsubsection{选修主题详细大纲}
\begin{itemize}
    \item \textbf{命题与一阶逻辑基础回顾}：规划状态的形式化表示语言。
    \item \textbf{规划算子与世界状态规范}：STRIPS 算子体系（前置先决条件 Preconditions、添加列表 Add List、删除列表 Delete List）。
    \item \textbf{全序规划}：状态空间前向搜索与目标后向搜索。
    \item \textbf{偏序规划}：因果链概念、动作定序约束与因果威胁消解。
    \item \textbf{规划图与 GraphPlan 算法}：命题层、动作层、互斥约束传播与逆向无冲突解抽取。
    \item \textbf{分层任务网络规划}：复合任务分解、非基本动作演化。
    \item \textbf{现代标准化规划描述语言}：规划领域定义语言 (PDDL)规范。
    \item \textbf{多智能体协同规划}：分布式动作协同与资源互斥。
    \item \textbf{基于马尔可夫决策过程的规划}。
    \item \textbf{闭环规划与动态重规划体系}：
    \begin{itemize}
        \item 条件规划（Conditional Planning / 感知-动作协同交织）；
        \item 连续规划；
        \item 概率规划。
    \end{itemize}
\end{itemize}

\subsubsection{示范性学习成果与能力目标}
\begin{enumerate}
    \item 针对给定的实际工程规划问题（如物流仓储包裹搬运或无人机多航点巡航），准确\textbf{构建}其状态表示、目标规范以及包含先决条件与动作效果的规划算子集。
    \item 熟练运用 PDDL 语言\textbf{形式化编写}规划领域的领域文件与问题实例文件，并调用现代化规划求解器（如 Fast Downward）求解出可行或最优动作计划。
    \item 给定初始状态、目标状态与算子集，规范绘制偏序规划因果图，准确识别因果威胁冲突，并通过添加定序因果约束\textbf{成功消解所有冲突}。
    \item 针对具体的规划问题实例，严密\textbf{构建完整的 GraphPlan 规划图}，系统推演命题层与动作层的互斥约束关系，并演示如何逆向抽取有效执行计划。
\end{enumerate}

\subsection{智能体系统 [选修]}
\textbf{【学时建议】}：拓展选修主题（与 人机交互/Collaboration and Communication 交叉）。

本单元系统讲授自主智能体与多智能体系统 (MAS)的核心架构、博弈协作机制与群体智能算法。

\subsubsection{选修主题详细大纲}
\begin{itemize}
    \item \textbf{智能体经典架构体系}：简单反应式架构、分层包容架构、认知架构（Cognitive Architectures，如 BDI、ACT-R、SOAR）。
    \item \textbf{智能体形式化理论}：信念-愿望-意图（BDI）多模态逻辑系统、智能体状态的形式化定义。
    \item \textbf{理性决策与博弈论}：纳什均衡、占优策略、零和博弈与非零和博弈。
    \item \textbf{决策论智能体与马尔可夫决策过程 (MDP)}。
    \item \textbf{软件智能体与智能助手}：个人虚拟助手、信息检索与过滤智能体。
    \item \textbf{协作智能体与信息搜集智能体}。
    \item \textbf{拟真与情感智能体}：虚拟合成角色、智能体情绪建模与拟人化交互。
    \item \textbf{学习型智能体体系}：学习元件、执行元件、评判元件与问题生成器。
    \item \textbf{多智能体系统}：
    \begin{itemize}
        \item 协作多智能体系统与智能体团队编队协作；
        \item 竞争性智能体交互机制：拍卖机制（Auctions，如英式拍卖、荷式拍卖、维克里第二价格拍卖）、投票协议与社会选择理论；
        \item 群体智能系统与生物启发式模型（Swarm Systems，如蚁群算法、粒子群优化算法 PSO）；
        \item 多智能体强化学习 (MARL)。
    \end{itemize}
    \item \textbf{人-智能体交互与协作}：混合主动性协同决策。
    \item \textbf{智能体通信方法学}：知识查询与操作语言（KQML）、FIPA 智能体通信语言（ACL）、合同网协议。
    \item \textbf{工程实践挑战与应用落地}：自动化金融高频交易智能体、智能分布式供应链协调调度。
\end{itemize}

\subsubsection{示范性学习成果与能力目标}
\begin{enumerate}
    \item 深入\textbf{阐释并对比}主流智能体架构（简单反射式、基于模型、基于目标、分层包容及认知架构）的控制机理与环境适用性。
    \item 深入\textbf{描述}智能体理论在纯软件自动化脚本、智能个人数字助手以及游戏可信拟真虚拟化身中的典型工程落地路径。
    \item 详尽\textbf{阐明}自主学习型智能体的四大核心构件及其相互之间的演进迭代机制。
    \item 结合生动具体的现实工程案例（如分布式传感器网络协作或分布式电力市场撮合交易），详尽\textbf{论证}多智能体系统如何通过交互协议与黑板机制实现高效协同。
    \item 熟练运用成熟的认知计算架构（如 ACT-R 或 SOAR），针对特定的复杂推理或人机交互任务，独立\textbf{设计并构建}一个具备长期记忆与技能学习能力的智能体原型系统。
\end{enumerate}

\subsection{自然语言处理 [选修]}
\textbf{【学时建议】}：拓展选修主题（与 人机交互/New Interactive Technologies 交叉）。

本单元系统讲授计算机处理、理解与生成人类自然语言的理论体系，涵盖经典文法解析、统计 NLP、信息检索以及基于 Transformer 与大语言模型（LLM）的现代深度学习前沿。

\subsubsection{选修主题详细大纲}
\begin{itemize}
    \item \textbf{文法体系}：确定性文法与随机/统计文法。
    \item \textbf{句法分析算法}：
    \begin{itemize}
        \item 上下文无关文法与图表句法分析器（Chart Parsers，如 CYK 动态规划算法）；
        \item 概率上下文无关文法与带权 CYK 语法分析树推导。
    \end{itemize}
    \item \textbf{语义表示与语义理解}：
    \begin{itemize}
        \item 基于形式逻辑的知识表示与语义演算；
        \item 语义角色标注 (SRL)与深层语义依存分析；
        \item 时序逻辑语义表示；
        \item 对话理解中的信念、愿望与意图（BDI）语义映射。
    \end{itemize}
    \item \textbf{语料库驱动的统计 NLP 方法}：
    \begin{itemize}
        \item $N$-元文法语言模型（$N$-grams）与隐马尔可夫模型（HMM）；
        \item 概率平滑与回退策略（Backoff，如 Good-Turing 平滑、Kneser-Ney 平滑）；
        \item 核心应用实战：词性标注与形态学词法分析。
    \end{itemize}
    \item \textbf{信息检索与抽取}：
    \begin{itemize}
        \item 向量空间模型；
        \item 词频-逆文档频率（TF-IDF）权重度量与余弦相似度计算；
        \item 检索系统评估指标：精确率、召回率、平均准确率均值（MAP）；
        \item 信息抽取技术（命名实体识别 NER、实体关系抽取）；
        \item 机器翻译（统计机器翻译 SMT 与神经机器翻译 NMT）。
    \end{itemize}
    \item \textbf{文本分类与主题建模}：词袋模型与高维稀疏特征表示。
    \item \textbf{自然语言深度学习前沿}：
    \begin{itemize}
        \item 词向量表示（Word2Vec 连续词袋/Skip-gram 模型、GloVe 词嵌入）；
        \item 循环神经网络（RNN、LSTM、GRU）的时序语义建模；
        \item 变换器自注意力架构（Transformers / 多头自注意力 Multi-head Attention、位置编码 Positional Encoding）；
        \item 预训练语言模型演进（BERT 双向自编码器、GPT 自回归解码器、T5 统一序列到序列架构）；
        \item 多模态图文联合嵌入（Multi-modal Embeddings，如 CLIP 跨模态表征对齐）；
        \item 生成式大语言模型与提示词工程。
    \end{itemize}
\end{itemize}

\subsubsection{示范性学习成果与能力目标}
\begin{enumerate}
    \item 准确界定并\textbf{对比剖析}确定性文法与统计随机文法，结合自然语言中广泛存在的句法结构歧义，举例说明引入概率文法模型的重要性与必要性。
    \item 模拟、推导并\textbf{编程实现}经典确定性句法分析算法或概率 CYK 动态规划算法，实现输入自然语言句子的层次化语法树解析。
    \item 深入\textbf{识别并阐述}人类自然语言的多重歧义性、语境高度依赖性及隐含常识对计算机实现深层语义结构化表示所构成的核心理论与工程挑战。
    \item 详述在自然语言处理中采用标准化大型基准语料库的核心价值，并系统\textbf{列举}针对语法分析、命名实体识别、阅读理解等任务的主流标准语料集。
    \item 熟练\textbf{梳理并对比}现代信息检索、神经机器翻译与大规模文本自动分类的核心算法管线与技术范式。
    \item 从零编码\textbf{实现 TF-IDF 特征提取模块}，对大规模非结构化文本语料进行特征工程降维与数值向量化，并训练调用现成的机器学习分类器构建一个端到端的高性能文本分类系统。
\end{enumerate}

\subsection{机器人学 [选修]}
\textbf{【学时建议】}：拓展选修主题（与 HCI 及 专用平台/Industrial Platforms 交叉）。

本单元系统讲授自主机器人系统的核心软硬件原理，涵盖多源传感感知、空间运动学变换、定位建图（SLAM）、运动路径规划以及人机协作安全控制。

\subsubsection{选修主题详细大纲}
\begin{itemize}
    \item \textbf{机器人学全貌与演进}：核心科学问题、技术演进脉络与尖端前沿进展。
    \item \textbf{尖端机器人传感与执行系统硬件体系}：
    \begin{itemize}
        \item 多源传感器配置：激光雷达（LIDAR）、超声波声呐、单目/双目立体视觉相机、RGB-D 深度相机、仿生事件相机、麦克风阵列、六维力敏触觉传感器等；
        \item 机械执行机构：轮式差速/全向驱动底盘、多自由度连续机械臂、仿生多指灵巧手抓具；
        \item 传感原始信号预处理、标定与数字滤波。
    \end{itemize}
    \item \textbf{机器人控制架构演进}：慎思式控制 vs. 反应式控制、混合分层控制、布雷滕伯格车辆模型。
    \item \textbf{世界模型与物理环境表征}：占用网格地图、拓扑图、八叉树地图与三维语义场景图。
    \item \textbf{物理世界固有不确定性}：传感器测量噪声、机械传动间隙与打滑、非结构化环境干扰。
    \item \textbf{空间几何描述与坐标变换}：
    \begin{itemize}
        \item 二维与三维空间位姿表示；
        \item 齐次变换矩阵、旋转矩阵与四元数运算。
    \end{itemize}
    \item \textbf{构型空间与环境地图}：构型空间维度、障碍物在 C-空间的膨胀映射、工作空间向构型空间的转换。
    \item \textbf{不确定性传感数据融合与状态估计}：卡尔曼滤波、扩展卡尔曼滤波（EKF）、粒子滤波（Particle Filter / 蒙特卡洛定位 MCL）。
    \item \textbf{自主移动机器人关键导航技术}：
    \begin{itemize}
        \item 同时定位与建图 (SLAM)；
        \item 正向运动学 (FK)与逆向运动学 (IK)；
        \item 运动路径规划与轨迹平滑优化：基于采样的规划算法（概率路图法 PRM、快速探索随机树 RRT 及渐近最优算法 RRT*）、人工势场法；
        \item 关节伺服动力学建模与顺应性阻抗控制；
        \item 基于机器视觉的伺服闭环控制。
    \end{itemize}
    \item \textbf{多机器人协作与群体机器人}：多机通信拓扑、编队控制、分布式任务分配与拍卖协商。
    \item \textbf{人-机协作与共融交互}：
    \begin{itemize}
        \item 物理人机共享工作空间安全监控与防碰撞机制；
        \item 人类作业意图预测与行为轨迹推断；
        \item 显式与隐式多模态人机交互（语音指令、肢体手势、视线跟踪）；
        \item 人机交互信任度建模与评估。
    \end{itemize}
\end{itemize}

\subsubsection{示范性学习成果与能力目标}
\begin{infobox}[教学实现模式说明]
鉴于实体工业级机器人硬件较为昂贵，本单元的所有实践与考核成果均可依托高保真仿真环境（如 Gazebo、Isaac Sim、Webots）或低成本开源教育机器人平台完成。
\end{infobox}
\begin{enumerate}
    \item 系统梳理并\textbf{归纳列举}当今顶尖机器人系统的核心能力与物理局限性，详述各类核心车载传感器（LIDAR、深度视觉、IMU）的工作原理及其不可或缺的数据处理流水线。
    \item 针对特定的自动化作业任务，综合集成外设传感器、驱动执行器与底层控制软件，独立\textbf{搭建一套完整的软硬件闭环机器人系统}。
    \item 熟练运用慎思式、反应式或混合分层控制架构，为移动机器人\textbf{编写控制程序}以自主实现简单作业任务。
    \item 深入理解构型空间几何映射机制，并在高维构型空间中\textbf{实现基础的运动路径规划算法}（如 A*、RRT 或 Bug 导航算法）。
    \item 定量刻画常见工业级传感器与执行机构的物理不确定性与噪声漂移，并熟练\textbf{制定并实现}针对这些不确定性的算法滤波与状态估计补偿方案。
    \item 深入\textbf{辨析}机器人在其内部表示外部物理世界时的各类拓扑与几何表示方案（栅格地图、点云地图、拓扑图），客观评述各自的优势与不足。
    \item 详尽\textbf{对比剖析}在已知地图与完全未知探索环境下机器人自主导航的主流技术路线（全局拓扑导航、局部实时避障 DWA 等），剖析各自的适用环境与鲁棒性边界。
    \item 深入\textbf{阐释}至少一种协调多台自主机器人协同感知与同步动作以合力完成复杂搬运或搜救任务的分布式协同算法。
    \item 深入\textbf{对比}“多机协同”与“人机共融协作”在动力学安全规范、意图辨识与突发交互容错上的核心差异，并将差异根源追溯至人机物理与认知环境的不确定性本质。
\end{enumerate}

\subsection{感知与计算机视觉 [选修]}
\textbf{【学时建议】}：拓展选修主题。

本单元系统讲授计算机从图像、视频、音频及多模态物理传感器信号中抽取结构化知识的高级理论，涵盖传统图像几何处理、特征工程、音频信号解析与现代深度感知学习网络。

\subsubsection{选修主题详细大纲}
\begin{itemize}
    \item \textbf{计算机视觉成像与基础表示}：数字图像采集几何光学原理、空间采样与灰度量化、多通道矩阵表示、色彩空间（RGB、HSV、Lab）变换与图像底层属性。
    \item \textbf{图像底层特征工程与几何分析}：形状表示、边缘检测算子（Sobel、Canny）、角点检测（Harris、FAST）、局部尺度不变特征描述符（SIFT、SURF、ORB）。
    \item \textbf{图像分割技术}：自适应阈值分割（Otsu 算法）、基于边缘跟踪的分割、区域生长算法与图割。
    \item \textbf{运动分析与光流估计}：光流约束方程、Lucas-Kanade 稀疏光流与 Horn-Schunck 稠密光流算法。
    \item \textbf{生成式视觉前沿}：扩散概率模型、神经辐射场、三维高斯泼溅（3D Gaussian Splatting）。
    \item \textbf{音频与语音识别基础}：连续声学信号数字化采样、短时傅里叶变换（STFT）、梅尔频率倒谱系数（MFCC）、端到端语音识别声学模型。
    \item \textbf{触觉与本体感受}：触觉阵列传感器信号解析、机器人自身运动力矩感知。
    \item \textbf{其他多感官跨模态感知}（如电子鼻气味传感等）。
    \item \textbf{感知模式识别范式与模块化架构}（与 人工智能/机器学习 深度交叉）。
    \item \textbf{分类与目标检测评估指标体系}：Top-1/Top-5 准确率、交并比、平均精度均值。
    \item \textbf{统计模式识别技术}。
    \item \textbf{深度感知学习前沿架构}：现代卷积神经网络（ResNet、ConvNeXt）、视觉变换器、端到端目标检测网络（YOLO 系列、DETR）、实例分割与大模型分割范式（Mask R-CNN、Segment Anything Model, SAM）。
\end{itemize}

\subsubsection{示范性学习成果与能力目标}
\begin{enumerate}
    \item 深入\textbf{阐述}图像与目标识别在现代人工智能技术生态中的基石地位，并系统列举其在工业自动化质检、自动驾驶、智慧医疗影像筛查等关键领域的重大产业应用。
    \item 详尽\textbf{列举}至少三种图像分割的核心算法范式（阈值分割、边缘算子分割、区域生长/图割法），剖析各自的算法数学定义、计算瓶颈与适用场景。
    \item \textbf{编码实现}基于几何轮廓边缘描述或二值区域特征的二维工业零件几何匹配与目标识别算法。
    \item 严格\textbf{辨析}声音事件检测、连续语音转写与特定说话人声纹确认在感知目标上的本质差异，指明原始音频时序信号在上述三类任务中预处理与特征抽取的根本不同。
    \item 举出至少两个在跨感官模态之间实现感知数据域转换的工程范例（例如将指尖高密度微触觉力敏阵列信号映射转化为单通道二维灰度压力分布图像进行视觉化卷积解析）。
    \item 在真实物理场景数据上独立\textbf{编程实现特征提取算子}（如实现高精度 Canny 边缘算子、Harris 亚像素角点检测器，或提取音频短时切片的快速傅里叶变换 FFT 频谱特征向量）。
    \item 编写算法将底层原始基元特征有效组合凝聚为高阶感知结构（例如将离散像素边缘连接拟合为凸多边形或平滑样条曲线，或将音频声学特征向量聚类映射为候选音素假设）。
    \item \textbf{编程实现}一个完备的模式分类器，将输入的连续感知特征向量正确划分至目标类别空间，并运用混淆矩阵对其泛化分类性能进行客观定量评估。
    \item 运用严格的控制变量对照实验，定量\textbf{评估}底层特征提取算法相较于其他备选特征提取机制对最终端到端模式识别准确率的独立贡献度与鲁棒性表现。
    \item 深入\textbf{阐释}至少三种经典感知分类算法的数学原理、先验假设前提、技术优势以及在复杂高维非凸地形下的性能短板。
    \item 针对计算机视觉中的复杂感知任务（如大规模通用物体识别、开集检测或复杂街景多目标语义分割），设计、微调并部署一个基于现代深度神经网络（如先进 CNN 或 ViT 架构）的工业级感知解决方案，并完成高标准的定量评估与算力推理时延调优。
\end{enumerate}

\section{专业素养与思维倾向}

在人工智能领域的理论研究与工程实践中，学生应在专业思维与职业素养方面重点培育以下三种核心倾向：

\begin{itemize}
    \item \textbf{严谨细致}：在实现人工智能和机器学习算法时，必须对细微之处保持高度敏锐与专注。数据分布的细微偏移、微小的超参数变动、张量维度的不对齐或概率平滑因子的微调，均可能引发灾难性的模型震荡或隐蔽的系统漏洞。学生必须具备对数学推导、数据管道及底层代码实现精益求精、一丝不苟的科学态度。
    \item \textbf{坚韧执着}：人工智能技术往往运行在部分可观察、高噪声与非平稳的复杂开放物理环境中，优化算法在多轮梯度迭代中极易遭遇误差级联放大、梯度消失或陷入鞍点。要使一个 AI 系统具备确定性、高鲁棒的工程落地表现，通常需要经历漫长的参数调优、假设验证、失败归因与反复试错。这种极具挑战性的技术特性要求学生在面对复杂阻碍与模型失效时，具备强大的心理韧性与持之以恒的攻关毅力。
    \item \textbf{责任与伦理担当}：人工智能的应用已深远渗透至社会生产与日常生活的各个角落，其输出结果直接影响微观个体的合法权益、人身财产安全以及宏观社会的公平与稳定。这要求学生深刻洞悉 AI 技术带来的跨学科社会影响，始终秉持对人类福祉高度负责的伦理底线，审慎决策何时何地适宜应用 AI，并自觉落实算法透明度、数据隐私保护与无歧视的伦理规范。
\end{itemize}

\section{数学基础要求}

人工智能作为一门理论深厚、高度跨学科的硬核计算机科学分支，要求学生具备扎实、严密的数理科学底座：

\subsection{必修数学基础}
\begin{itemize}
    \item \textbf{离散数学}：
    \begin{itemize}
        \item 集合论（集合代数、笛卡尔积、幂集）、二元关系（等价关系、偏序关系）、映射与函数理论；
        \item 命题逻辑与一阶谓词逻辑语法与语义、逻辑推理规则、形式化数学归纳法与反证法。
    \end{itemize}
    \item \textbf{线性代数}：
    \begin{itemize}
        \item 矩阵代数系统：矩阵乘法、转置、行列式、伴随矩阵、矩阵求逆、秩与线性方程组求解；
        \item 向量空间理论：线性无关性、正交基底、子空间投影、特征值与特征向量分解、奇异值分解（SVD）。
    \end{itemize}
    \item \textbf{概率论与数理统计}：
    \begin{itemize}
        \item 基础概率公理系统、条件概率计算、全概率公式、事件独立性；
        \item 贝叶斯定理及其在逆向推断中的严谨应用；
        \item 随机变量数学期望（均值）、方差、协方差、高阶矩；
        \item 经典概率分布族（高斯正态分布、伯努利分布、二项分布、泊松分布、多项分布）；
        \item 描述性统计量与推断统计学基础（假设检验、$p$-值检验、置信区间估计）；
        \item \textbf{【教学特别强调】}：上述所有概率统计理论必须紧密围绕真实工程决策场景、运用真实世界“脏数据”开展案例教学，坚决摒弃脱离现实工程本质的纯人工玩具例题。
    \end{itemize}
\end{itemize}

\subsection{进阶与推荐数学基础}
\begin{itemize}
    \item \textbf{基于微积分的高等概率与连续统计}：多维连续型随机变量联合概率密度函数、多重微积分变换、概率生成函数与特征函数。
    \item \textbf{高级推断与统计实验设计}：科学实验设计方法学、数据重采样与非参数方法（Bootstrap 自助抽样法、Jackknife 刀切法、排列检验）。
    \item \textbf{数值优化理论}：凸集与凸函数理论、多变量连续函数无约束梯度优化、拉格朗日乘子法与 KKT 最优性条件、对偶理论、牛顿法与拟牛顿法（L-BFGS）。
\end{itemize}

\section{交叉主题与跨学科共享概念}

\subsection{跨领域共享概念}
\begin{itemize}
    \item \textbf{搜索算法}：人工智能/搜索算法与问题求解 单元与 算法/Algorithms and Complexity（算法基础与复杂性）中的高级图算法和搜索策略深度融合。
    \item \textbf{数据管理与数据工程}：机器学习全流程中的高维数据治理、特征仓库、分布式数据读取与 数据管理/Data Management 及数据科学基础设施紧密贯通。
    \item \textbf{伦理与社会规范}：人工智能/人工智能应用与社会影响 单元与 社会伦理/Society, Ethics and Professionalism（社会、伦理与专业素养）构建了统一的科技治理视角。
    \item \textbf{机器人学与嵌入式硬件}：人工智能/机器人学 单元与 专用平台/Specialized Platform Development（专用平台开发中的机器人工业控制平台）构成软硬件一体化的协同映射。
\end{itemize}

\subsection{贯穿式核心维度}
\begin{itemize}
    \item \textbf{计算效率}：贯穿于高维状态搜索的剪枝加速、图模型推断的因果解耦，以及大规模深度神经网络的分布式算力并行优化中。
    \item \textbf{伦理道德}：贯穿于从数据采集标注、特征工程设计，到模型训练偏见控制、可解释性审查及司法/医疗落地全生命周期的伦理约束。
    \item \textbf{系统建模}：贯穿于将开放非结构化现实业务难题抽象转化为状态图搜索、确定性逻辑命题、概率图模型以及损失函数优化数学模型的建模全过程。
    \item \textbf{工程编程}：强调将复杂的数学优化方程转化为高内聚、低耦合、高吞吐的工业级高质量生产代码实现。
\end{itemize}

\section{课程设计与开课规划建议}

联合工作组针对不同培养侧重，提出三套典型的人工智能课程体系设计方案：

\subsection{方案一：《人工智能基础导论》（38 学时综合基石课程）}
面向全校或计算机大类专业通识进阶设置，旨在让学生系统建立对人工智能各核心领域的全景认知与扎实求解能力：
\begin{itemize}
    \item \textbf{学时构成}：
    \begin{itemize}
        \item 人工智能/人工智能基础议题（基础议题）：4 学时
        \item 人工智能/搜索算法与问题求解（搜索算法与问题求解）：9 学时
        \item 人工智能/知识表示与推理基础（基础知识表示与推理）：4 学时
        \item 人工智能/Probabilistic Knowledge Representation and Reasoning（概率表示与推理）：5 学时
        \item 人工智能/机器学习（机器学习）：12 学时
        \item 人工智能/人工智能应用与社会影响（应用与社会影响）：4 学时（强调整合贯穿全课程各讲）
    \end{itemize}
    \item \textbf{先修课程要求}：CS2（数据结构与面向对象编程）、离散数学、概率论基础。
    \item \textbf{能力目标陈述}：完成本课程的学生能够系统掌握人工智能各大核心子领域的基本范式与方法论，具备在各方向独立理解、开发并应用代表性算法的核心能力；能够运用状态搜索技术、基础贝叶斯推理与经典机器学习方法求解复杂的现实工程问题；能够全面理解 AI 技术的广泛应用潜力，并能对其引发的法律合规、算法偏见与社会伦理挑战进行理性的审视与权衡。
\end{itemize}

\subsection{方案二：《机器学习原理与实战》（44 学时深度算法课程）}
面向数据科学、智能科学与技术或研究型计算机专业方向专设，聚焦于数据驱动的学习系统理论与工业级实战：
\begin{itemize}
    \item \textbf{学时构成}：
    \begin{itemize}
        \item 人工智能/机器学习（机器学习与深度网络前沿）：32 学时
        \item 人工智能/知识表示与推理基础（知识表示与贝叶斯推断）：4 学时
        \item 人工智能/自然语言处理（自然语言处理精选专题：TF-IDF、词袋模型、文本分类与 Transformer 初步）：4 学时
        \item 人工智能/人工智能应用与社会影响（应用与伦理公平审计）：4 学时（贯穿全生命周期实验教学）
    \end{itemize}
    \item \textbf{先修课程要求}：CS2、离散数学、概率论、数理统计、线性代数（强烈推荐）。
    \item \textbf{能力目标陈述}：完成本课程的学生能够透彻理解、自主设计、独立开发并部署先进的监督学习与强化学习模型；能够针对特定工程问题准确选择最适宜的算法范式，严谨执行无泄漏的数据预处理与特征工程，设计规范科学的评估基准，并深入解释产出模型的决策逻辑与其内在缺陷；能够敏锐识别并消除有偏数据集及各类系统误差来源，在将机器学习应用于实际工程问题时严格践行符合伦理道德与社会责任的标准。
\end{itemize}

\subsection{方案三：《自主机器人系统与智能控制》（42 学时机电一体化软硬件课程）}
面向智能装备、物联网与具身智能领域专设，强调软硬件深度交织的物理交互与环境闭环导航：
\begin{itemize}
    \item \textbf{学时构成}：
    \begin{itemize}
        \item 人工智能/机器人学（机器人学核心理论与实践）：25 学时
        \item 专用平台/机器人平台（专用机器人硬件平台、物理约束与嵌入式软件架构）：4 学时
        \item 人工智能/搜索算法与问题求解（机器人路径搜索与避障规划，如 A*、RRT 算法）：4 学时
        \item 人工智能/机器学习（机器学习在机器人感知与视觉目标识别中的应用）：6 学时
        \item 人工智能/人工智能应用与社会影响（自主机器人系统社会伦理与安全影响）：3 学时（融入全流程安全论证）
    \end{itemize}
    \item \textbf{先修课程要求}：CS2、线性代数与空间解析几何。
    \item \textbf{能力目标陈述}：完成本课程的学生能够系统掌握先进机器人学技术体系；熟练运用多源传感器感知物理外部世界，基于提取的环境特征与高精度地图实现鲁棒的自主实时定位（SLAM），并在简单及中等复杂度的现实场景中完成运动路径规划与自主移动闭环导航；能够熟练运用基础计算机视觉处理算法、机械臂正逆运动学解析求解技术及空间轨迹优化方法。
\end{itemize}

\section{真实能力任务规范}

真实能力规范是 CS2023 引入的重大教育创新，通过高还原度的真实工程挑战，综合考量学生的知识融通、高阶认知与专业素养。

\begin{taskbox}[任务 1：销售库存智能优化决策工具]
\textbf{【任务情境】}：某全国性连锁商超或跨国零售企业的区域销售团队面临数万种 SKU 商品在多门店与中心仓之间的动态调配难题，需要在满足复杂货架库容、冷链保鲜时效及运输履约成本多重严苛物理制约的前提下，实现库存调配方案的全局动态优化。\\
\textbf{【能力陈述】}：将该业务库存调度问题规范抽象为状态空间图搜索模型，准确界定显式与隐式业务边界约束，设计具有严密理论可采纳性的高品质启发函数，编码实现智能化图搜索算法（如 A* 搜索结合约束满足 CSP 求解器），高效求解出满足各项严苛约束的最优或近优库存调配执行方案。\\
\textbf{【能力领域】}：应用 \quad \textbf{【能力单元】}：系统设计 / 软件开发 / 性能评估\\
\textbf{【所需知识单元】}：人工智能/搜索算法与问题求解\\
\textbf{【认知与技能要求】}：\textbf{应用 / 开发} \quad \textbf{【核心要求】}：计算机核心
\end{taskbox}

\begin{taskbox}[任务 2：不确定环境下双人博弈策略智能体]
\textbf{【任务情境】}：设计一个在复杂双人零和对抗博弈（如包含随机抽卡或掷骰子机会事件的棋盘游戏，或带有突发价格波动与信息不完全的战略采购博弈）中能够自主做出全局最优决策的博弈智能体。\\
\textbf{【能力陈述】}：熟练构建带有 Alpha-Beta 剪枝的极小化极大算法（Minimax with $\alpha$-$\beta$ Pruning），引入处理随机不确定性概率动作的期望机会节点，并在指定搜索深度处设计合理的启发式局面估值函数，开发并部署一个能在不确定对抗环境下具备高胜率对抗水平的智能博弈决策引擎。\\
\textbf{【能力领域】}：应用 \quad \textbf{【能力单元】}：需求分析 / 架构设计 / 软件开发 / 仿真测试 / 生产部署\\
\textbf{【所需知识单元】}：人工智能/搜索算法与问题求解、人工智能/知识表示与推理基础\\
\textbf{【认知与技能要求】}：\textbf{应用 / 开发} \quad \textbf{【核心要求】}：领域核心
\end{taskbox}

\begin{taskbox}[任务 3：基于多源体检数据的疾病患病概率贝叶斯预测系统]
\textbf{【任务情境】}：智慧医疗系统需要基于病人体检报告中呈现的非完美化验指标、不完全既往病史及流行病学先验统计，对特定高危慢性病或急性并发症的发病风险给出具备统计置信度的概率推断。\\
\textbf{【能力陈述】}：将现实临床诊断中的多维不确定性因果关联建模为概率贝叶斯推理问题，熟练应用贝叶斯法则计算后验概率，编码实现精确或采样概率推断算法，根据输入的临床观测证据实时推导出疾病发生率并给出具备统计说服力的合理解释。\\
\textbf{【能力领域】}：应用 \quad \textbf{【能力单元】}：系统开发 / 实验测试\\
\textbf{【所需知识单元】}：人工智能/搜索算法与问题求解、人工智能/逻辑表示与推理、数学统计/概率与统计\\
\textbf{【认知与技能要求】}：\textbf{应用 / 开发} \quad \textbf{【核心要求】}：计算机核心
\end{taskbox}

\begin{taskbox}[任务 4：现代复杂物流调度与分拣规划求解器]
\textbf{【任务情境】}：大型现代化智能物流枢纽面临海量多目的地包裹的自动化交叉分拣、机械臂连续抓取位姿衔接与无人引导搬运车（AGV）无碰撞路径协同的时空连续规划挑战。\\
\textbf{【能力陈述】}：将该复杂物流调度任务形式化规范为规划初始状态、目标状态以及包含先决条件与动作效果的规划算子集，运用 PDDL 语言建模并采用偏序规划或分层任务网络规划算法，求解出无资源死锁且全局搬运效率最优的动作序列。\\
\textbf{【能力领域】}：应用 \quad \textbf{【能力单元】}：系统设计 / 理论阐释\\
\textbf{【所需知识单元】}：人工智能/逻辑表示与推理、人工智能/规划系统\\
\textbf{【认知与技能要求】}：\textbf{解释 / 应用} \quad \textbf{【核心要求】}：拓展选修
\end{taskbox}

\begin{taskbox}[任务 5：表格化用户消费行为监督学习预测流水线]
\textbf{【任务情境】}：某大型零售或金融科技企业积累了海量高维表格化交易日志数据（如用户历史消费明细、APP 浏览频次、人口统计学标签等），需挖掘关键消费趋势并对高价值客户的未来流失概率进行精准预测。\\
\textbf{【能力陈述】}：熟练运用工业级机器学习库（如 Scikit-Learn、XGBoost），自主构建端到端的数据预处理流水线（缺失值工程补全、高维类别特征 One-Hot 与 Target 编码、特征无偏缩放），严格隔离数据集以防御数据泄漏，实施分层交叉验证与超参数网格搜索调优，构建出具备高泛化能力的监督学习分类模型并进行专业质量评估。\\
\textbf{【能力领域】}：应用 \quad \textbf{【能力单元】}：需求分析 / 架构设计 / 算法研发 / 效能评估 / 生产部署\\
\textbf{【所需知识单元】}：人工智能/机器学习、人工智能/人工智能应用与社会影响\\
\textbf{【认知与技能要求】}：\textbf{应用 / 开发} \quad \textbf{【核心要求】}：计算机核心
\end{taskbox}

\begin{taskbox}[任务 6：基于深度迁移微调的商品货架图像识别系统]
\textbf{【任务情境】}：无人零售便利店与智能结算货架需要通过高分辨率工业相机拍摄的杂乱货架照片，实时高精度识别上百种摆放各异、存在光照变化与局部遮挡的零售商品品类。\\
\textbf{【能力陈述】}：掌握现代深度学习与迁移学习技术体系，基于先进的预训练视觉大模型（如 ResNet 或 Vision Transformer Backbone），设计任务专用的特征分类头，利用标注业务样本完成模型微调，使用严谨的混淆矩阵、Top-1 准确率与 mAP 指标全面评估模型质量，并将模型工程化导出部署至边缘计算设备进行实时推理。\\
\textbf{【能力领域】}：应用 \quad \textbf{【能力单元】}：需求分析 / 架构设计 / 算法研发 / 系统测试 / 效能评估 / 生产部署\\
\textbf{【所需知识单元】}：人工智能/机器学习、人工智能/人工智能应用与社会影响\\
\textbf{【认知与技能要求】}：\textbf{应用 / 开发} \quad \textbf{【核心要求】}：领域核心
\end{taskbox}

\begin{taskbox}[任务 7：视频游戏中的深度强化学习智能对抗体]
\textbf{【任务情境】}：在复杂即时战略或动作类视频游戏中，开发团队需要构建一个具备高水平博弈意识、能够自主适应玩家战术流派演进的拟人化智能非玩家角色（AI 对手）。\\
\textbf{【能力陈述】}：深刻理解并应用现代深度强化学习 (DRL)算法框架（如 PPO、DQN 或 Actor-Critic 架构），将原始游戏画面像素与物理状态向量构建为环境状态，设计多目标奖励函数，完成策略网络的端到端训练与超参数收敛性调优，成功部署并在复杂视觉输入下展现出高胜率竞技水平。\\
\textbf{【能力领域】}：应用 \quad \textbf{【能力单元】}：需求定义 / 算法设计 / 策略开发 / 仿真评估 / 游戏部署\\
\textbf{【所需知识单元】}：人工智能/机器学习\\
\textbf{【认知与技能要求】}：\textbf{应用 / 开发} \quad \textbf{【核心要求】}：领域核心
\end{taskbox}

\begin{taskbox}[任务 8：社交媒体多模态/文本情绪与情感分析系统]
\textbf{【任务情境】}：数字舆情监测平台需要对主流社交媒体上每日产生的数百万条用户短文本发帖与图文评论进行实时多维情感极性（正面、负面、中性及具体细粒度情绪）精准识别与趋势预警。\\
\textbf{【能力陈述】}：构建非结构化文本处理流水线，执行文本清洗、分词、停用词过滤与现代特征提取（如 TF-IDF 向量化结合密集词嵌入，或微调预训练 Transformer），训练并调优监督分类器完成文本情感分类，输出包含置信度与情感极性权重的可视化分析报告。\\
\textbf{【能力领域】}：应用 \quad \textbf{【能力单元】}：需求分析 / 系统设计 / 算法研发 / 效能评估 / 生产部署\\
\textbf{【所需知识单元】}：人工智能/机器学习、人工智能/自然语言处理、人工智能/人工智能应用与社会影响\\
\textbf{【认知与技能要求】}：\textbf{应用 / 开发} \quad \textbf{【核心要求】}：领域核心
\end{taskbox}

\begin{taskbox}[任务 9：已部署机器学习系统的数据集偏差与算法偏见审计校正]
\textbf{【任务情境】}：某大型金融机构或公共服务部门在生产环境中运行的自动化信贷审批与人才简历筛选模型，被外部公众或监管机构质疑存在针对特定性别、年龄层或少数族裔群体的系统性偏见。\\
\textbf{【能力陈述】}：深刻理解、准确定位并量化审计生产机器学习模型中的各类偏见根源（历史数据代表性缺失偏差、特征标注偏差、算法损失惩罚倾斜及评测基准失真），撰写符合合规标准的专业审计报告，设计算法层面的缓解修正策略（如反事实再加权、对抗去偏训练、决策阈值多重校准），显著降低不公概率并向治理委员会做技术汇报。\\
\textbf{【能力领域】}：应用 \quad \textbf{【能力单元】}：沟通协作与专业审计\\
\textbf{【所需知识单元】}：人工智能/机器学习、人工智能/人工智能应用与社会影响\\
\textbf{【认知与技能要求】}：\textbf{理解与阐述} \quad \textbf{【核心要求】}：计算机核心
\end{taskbox}

\begin{taskbox}[任务 10：基于大语言模型（LLM）的一线智能客服系统设计与风控对齐]
\textbf{【任务情境】}：跨国电商或电信运营商拟引入基于商用大语言模型（LLM）的一线智能客服系统，实现 7$\times$24 小时自然交互，但系统必须严防事实幻觉误导用户、敏感数据泄露或被提示注入攻击。\\
\textbf{【能力陈述】}：深刻理解现代生成式深度大模型的内在工作机理与潜在失效模式（如事实性幻觉 Hallucination、逻辑脆弱性、越狱诱导脆弱性等），合理设计检索增强生成（RAG）管道与精细化提示工程，搭建安全围栏与人工兜底回退机制，在充分释放大模型高情商对话能力的同时严控业务安全与法律合规风险。\\
\textbf{【能力领域】}：应用 \quad \textbf{【能力单元】}：沟通论证与技术方案设计\\
\textbf{【所需知识单元】}：人工智能/人工智能应用与社会影响\\
\textbf{【认知与技能要求】}：\textbf{理解与阐述} \quad \textbf{【核心要求】}：计算机核心
\end{taskbox}

\section{工作组委员会名单}

\textbf{工作组主席}：
\begin{itemize}
    \item Eric Eaton，宾夕法尼亚大学 (USA)
\end{itemize}

\textbf{工作组委员}：
\begin{itemize}
    \item Zachary Dodds，哈维穆德学院 (USA)
    \item Susan L. Epstein，纽约市立大学亨特学院与研究生中心（Hunter College \& The Graduate Center of CUNY, New York, NY, USA）
    \item Laura Hiatt，美国海军研究实验室 (USA)
    \item Amruth N. Kumar，新泽西拉马波学院 (USA)
    \item Peter Norvig，谷歌公司 / 斯坦福大学
    \item Meinolf Sellmann，通用电气全球研发中心 (USA)
    \item Reid Simmons，卡耐基梅隆大学 (USA)
\end{itemize}

\textbf{贡献专家}：
\begin{itemize}
    \item Claudia Schulz，汤森路透
\end{itemize}

\clearpage
\onecolumn
\section{附录：核心主题与技能要求等级对照表}

本附录完整汇总了 CS2023 人工智能知识领域中所有计算机核心与领域核心知识点、具体主题内容、布鲁姆认知技能达标要求等级及推荐核心学时分配：

\begingroup
\zihao{-5}
\renewcommand{\arraystretch}{0.92}
\setlength{\tabcolsep}{3pt}
\noindent\begin{tabularx}{\textwidth}{>{\raggedright\arraybackslash}p{0.14\textwidth} Y >{\centering\arraybackslash}p{0.15\textwidth} >{\centering\arraybackslash}p{0.09\textwidth} >{\centering\arraybackslash}p{0.06\textwidth}}
\toprule
\textbf{知识单元} & \textbf{核心知识点与具体主题} & \textbf{认知等级} & \textbf{分类} & \textbf{学时} \\
\midrule
\multirow{4}{*}{\makecell[l]{人工智能/基础议题}} 
& AI 问题概述、近期重大成功应用案例；智能体定义与典型范式（反应式、深思式）；什么是智能行为；图灵测试及其理论缺陷；多模态输入与输出；智能行为的计算模拟；理性推理 vs. 非理性/类人推理 & 解释 & 计算机核心 & \multirow{3}{*}{2} \\
\cline{2-4}
& 问题特征分类（完全/部分可观察、单/多智能体、确定/随机、静态/动态、离散/连续）；智能体本质与自主性特征（完全/半/混合自主，反射/目标/效用，不确定性与不完全信息决策，感知与环境交互核心地位，基于学习的智能体，具身智能体三要素：传感器、动力学、执行器） & 评估 & 计算机核心 & \\
\cline{2-4}
& 人工智能产业应用、爆发式增长与多维深远影响（经济、社会、法律与伦理治理） & 解释 & 计算机核心 & \\
\cline{2-5}
& 结合复杂现实工程实例对环境与问题特征的深度剖析；结合工业级案例对各类智能体架构与行为特性的深度分析；AI 发展趋势与宏观经济、社会、伦理冲击的深入研讨 & 评估 & 领域核心 & 1 \\
\midrule
\multirow{5}{*}{\makecell[l]{人工智能/搜索算法\\与问题求解}} 
& 状态空间表示（状态、目标与动作算子规范，假设空间状态因子化，基于图搜索的问题求解，图的隐式动态按需展开） & 解释 & 计算机核心 & \multirow{4}{*}{\makecell[c]{5\\{\footnotesize (含 3h}\\{\footnotesize 共享 AL)}}} \\
\cline{2-4}
& 无信息图搜索求解（广度优先搜索 BFS、深度优先搜索 DFS 及其迭代加深 IDS、一致代价搜索 UCS） & 开发/应用 & 计算机核心 & \\
\cline{2-4}
& 启发式图搜索求解（启发函数构建与可采纳性/一致性、爬山法、局部极小值与搜索地形、贪婪最佳优先搜索、A* 搜索） & 开发/应用 & 计算机核心 & \\
\cline{2-4}
& 图搜索算法的时间复杂度与空间复杂度理论推导、完备性与最优性严密评估 & 评估 & 计算机核心 & \\
\cline{2-5}
& 双向搜索、束搜索；双人对抗博弈（极小化极大搜索 Minimax、Alpha-Beta 剪枝技术、深度截断与启发式静态估值）；A* 算法的高效工程优化实现 & 开发/应用 & 领域核心 & 4 \\
\midrule
\multirow{3}{*}{\makecell[l]{人工智能/基础知识\\表示与推理}} 
& 知识表示类型（符号与逻辑表示、从自然语言抽象形式逻辑、学习驱动的亚符号连续向量嵌入、概率图模型初阶如朴素贝叶斯与贝叶斯网） & 解释 & 计算机核心 & \multirow{2}{*}{\makecell[c]{2\\{\footnotesize (含 1h}\\{\footnotesize 共享 MSF)}}} \\
\cline{2-4}
& 概率推理回顾、贝叶斯定理（与 MSF 离散概率交叉）；贝叶斯推断与证据置信度更新计算 & 应用 & 计算机核心 & \\
\cline{2-5}
& 随机变量与概率分布、概率公理体系、概率推断机制、贝叶斯法则数学推导、复杂场景下的贝叶斯推理与证据更新、边际独立性与条件独立性、效用理论与最大期望效用决策 & 应用 & 领域核心 & 2 \\
\midrule
\multirow{5}{*}{\makecell[l]{人工智能/机器学习}} 
& 机器学习主流任务范式（监督学习：分类与回归；强化学习；无监督学习：聚类）；理论哲学（没有免费的午餐定理、不可判定性、系统误差来源） & 应用/开发/评估 & 计算机核心 & \multirow{3}{*}{4} \\
\cline{2-4}
& 简明统计监督学习（朴素贝叶斯、决策树原理与调用）；过拟合与复杂度控制（正则化、剪枝直观理解，偏差-方差权衡）；真实数据工程（清洗、缺失值填补/标记、特征编码、归一化/标准化、真实脏数据处理）；特征表示（假设空间、基特征扩展、表示学习）；模型严谨评测（训练/验证/测试集隔离、分类评估指标、保留测试集性能估计、验证集超参数调优、模型局限洞察）；基础神经网络（前向与反向传播直观理解） & 应用/开发/评估 & 计算机核心 & \\
\cline{2-4}
& 机器学习工程伦理（真实场景案例、数据集偏见、算法偏见、评估偏见识别） & 解释/评估 & 计算机核心 & \\
\cline{2-5}
& 形式化优化建模（最小二乘线性回归、逻辑回归、损失函数、梯度下降法、过拟合正则化数学推导）；模型集成（加权多数表决）；深度学习（深度前馈网络、CNN 局部感受野/权重共享直观机制、深度网络学到特征表征的可视化）；高级性能评估（分类高级指标、回归指标、混淆矩阵、交叉验证、超参数网格/随机搜索）；强化学习全景概览；重大跨行业工程落地剖析 & 应用/开发/评估 & 领域核心 & \multirow{2}{*}{4} \\
\cline{2-4}
& 机器学习伦理深化（真实案例研究、差分隐私与数据隐私保护、算法公平性度量） & 解释/评估 & 领域核心 & \\
\midrule
\multirow{4}{*}{\makecell[l]{人工智能/应用与\\社会影响}} 
& 跨学科重大领域实战应用（智慧医疗、可持续发展、社交网络、金融经济、教育、机器人等，抽象为 AI 求解问题、数据可用性与清洗、数据/算法/评估偏见防范） & 解释/评估 & 计算机核心 & \multirow{2}{*}{2} \\
\cline{2-4}
& 已部署前沿深度生成式模型（深度图像扩散生成模型、大语言模型 LLM 的机制、落地赋能与事实幻觉/逻辑脆弱性失效短板）；AI 深刻社会影响（伦理、算法公平性、信任度与决策可解释性） & 解释/评估 & 计算机核心 & \\
\cline{2-5}
& 深入剖析一个或多个额外重大应用垂直领域；大语言模型与深度生成模型的工业实战前沿与高安全场景落地边界；社会治理议题（伦理责任、防歧视公平治理、可信度、因果可解释性）的纵深批判性研讨 & 解释/评估 & 领域核心 & 2 \\
\bottomrule
\end{tabularx}
\endgroup
\clearpage
\twocolumn
